\documentclass[11pt]{article}

\usepackage[margin=1in]{geometry}
\usepackage[T1]{fontenc}
\usepackage{lmodern}
\usepackage{microtype}
\usepackage{amsmath,amssymb,amsthm,mathtools}
\usepackage{graphicx}
\usepackage[hidelinks]{hyperref}

\newtheorem{theorem}{Theorem}[section]
\newtheorem{proposition}[theorem]{Proposition}
\newtheorem{lemma}[theorem]{Lemma}
\newtheorem{corollary}[theorem]{Corollary}
\theoremstyle{remark}
\newtheorem{remark}[theorem]{Remark}

\newcommand{\R}{\mathbb{R}}
\newcommand{\E}{\mathbb{E}}
\newcommand{\Id}{I}
\newcommand{\Rad}{\operatorname{Rad}}
\newcommand{\spann}{\operatorname{span}}
\newcommand{\codim}{\operatorname{codim}}
\newcommand{\restr}[2]{#1\!\restriction_{#2}}

\title{A Negative Answer to Ostrovskii's\\Projection-Factorization Problem}
\author{}
\date{June 19, 2026}

\begin{document}
\maketitle

\begin{abstract}
Ostrovskii asked whether there is a universal constant controlling a projection
onto an intermediate finite-dimensional space while requiring its composition
with a second projection to be a minimal projection onto a smaller space.  We
first verify an elementary correction argument that gives the dimension-dependent
upper bound $1+2\sqrt d$, where $d$ is the dimension of the smaller space.  We
then give a counterexample to the universal-constant question.  More precisely,
for every $d\geq 2$ we construct a finite-dimensional triple
$X_{1,d}\subset X_{2,d}\subset X_{3,d}$ with $\dim X_{1,d}=d$ and
$\lambda(X_{2,d},X_{3,d})\leq 2$ such that every pair of projections
$P_1:X_{2,d}\to X_{1,d}$ and $P_2:X_{3,d}\to X_{2,d}$ satisfying
$\|P_1P_2\|=\lambda(X_{1,d},X_{3,d})$ must obey
\[
        \|P_2\|\geq \sqrt{d/6}-2.
\]
Consequently no universal constant exists.  The only non-elementary input is
the known unique minimality of the classical Rademacher projection in
$L_p[0,1]$, $1<p<\infty$.
\end{abstract}

\section*{Source problem screenshot}

The source question is Problem~3 in M.~I. Ostrovskii,
\emph{Compositions of projections in Banach spaces and relations between
approximation properties}, arXiv:0811.1763.

\begin{center}
\includegraphics[width=\linewidth]{figures/open_problem_crop.png}
\end{center}

\section*{Status and idea}

This packet combines two compatible results.  First, the correction formula in
Section~\ref{sec:correction} proves a positive dimension-dependent theorem:
Problem~3 is always solvable with constant $C_d=1+2\sqrt d$ when
$d=\dim X_1$.  Second, Theorem~\ref{thm:main} gives a counterexample family to
the dimension-free version asked in the source paper.  Together they show that
the correct worst-case dependence on $\dim X_1$ is of order $\sqrt d$, up to
absolute constants.

The idea of the negative answer is to start with a large-norm projection that
is uniquely minimal.  A hyperplane containing its range then forces any
factorization whose product remains minimal to pass through a projection with
large norm.  The required uniquely minimal projections are the finite
Rademacher projections in $L_{2d}(\{-1,1\}^d)$, using the known
Lewicki--Skrzypek uniqueness theorem for the classical Rademacher projection.

\section{The problem and the main result}

For a finite-dimensional subspace $Y$ of a Banach space $X$, write
\[
 \lambda(Y,X)=\inf\{\|Q\|: Q:X\to Y\text{ is a projection}\}.
\]
Ostrovskii's Problem~3 \cite{Ostrovskii} asks whether there is a constant $C<\infty$ such that,
for every triple $X_1\subset X_2\subset X_3$ with $X_1$ and $X_2$
finite-dimensional, one can find projections
\[
 P_1:X_2\to X_1,\qquad P_2:X_3\to X_2
\]
for which
\[
 \|P_2\|\leq C\lambda(X_2,X_3)
 \quad\text{and}\quad
 \|P_1P_2\|=\lambda(X_1,X_3).
\]
All spaces in the counterexample below are real.  This is enough to disprove
the universal statement.

\begin{theorem}[Negative answer]\label{thm:main}
For every integer $d\geq 2$ there are finite-dimensional real Banach spaces
\[
 X_{1,d}\subset X_{2,d}\subset X_{3,d},
 \qquad \dim X_{1,d}=d,
\]
such that $\lambda(X_{2,d},X_{3,d})\leq 2$ and every pair of projections
$P_1:X_{2,d}\to X_{1,d}$ and $P_2:X_{3,d}\to X_{2,d}$ satisfying
\[
       \|P_1P_2\|=\lambda(X_{1,d},X_{3,d})
\]
necessarily satisfies
\[
       \|P_2\|\geq \sqrt{d/6}-2.
\]
In particular,
\[
 \frac{\|P_2\|}{\lambda(X_{2,d},X_{3,d})}
 \geq \frac{\sqrt{d/6}-2}{2},
\]
so Ostrovskii's Problem~3 has a negative answer.
\end{theorem}

The construction has $\dim X_{3,d}=2^d$ and
$\dim X_{2,d}=2^d-1$.  The proof is given in Sections~\ref{sec:obstruction}
and~\ref{sec:rademacher}.

\section{Verification of the dimension-dependent correction argument}
\label{sec:correction}

We first record the partial result from the supplied solution packet.  The
argument is correct; in fact it gives the slightly sharper unfactored estimate
in \eqref{eq:unfactored} below.

\begin{proposition}[Correction formula]\label{prop:correction}
Let $X_1\subset X_2\subset X_3$, where $X_1$ and $X_2$ are
finite-dimensional.  There exist projections $P_1:X_2\to X_1$ and
$P_2:X_3\to X_2$ such that
\[
 \|P_1P_2\|=\lambda(X_1,X_3)
\]
and
\begin{equation}\label{eq:unfactored}
 \|P_2\|\leq \lambda(X_2,X_3)+\lambda(X_1,X_3)
       +\lambda(X_1,X_3)\lambda(X_2,X_3).
\end{equation}
If $X_2\neq\{0\}$, then in particular
\[
 \|P_2\|\leq
 \bigl(1+2\lambda(X_1,X_3)\bigr)\lambda(X_2,X_3).
\]
\end{proposition}

\begin{proof}
Projection constants onto finite-dimensional ranges are attained.  Choose
minimal projections
\[
 P:X_3\to X_1,\qquad R:X_3\to X_2.
\]
Define
\[
 P_1=\restr{P}{X_2},\qquad P_2=R+P-PR.
\]
The range of $P_2$ is contained in $X_2$.  If $x\in X_2$, then $Rx=x$ and
therefore $P_2x=x+Px-Px=x$.  Thus $P_2$ is a projection onto $X_2$.
Moreover,
\[
 P_1P_2=P(R+P-PR)=PR+P-P^2R=P.
\]
Hence $P_1P_2$ is minimal onto $X_1$.  Finally,
\[
 \|P_2\|\leq \|R\|+\|P\|+\|P\|\,\|R\|,
\]
which is \eqref{eq:unfactored}.  Since $\lambda(X_2,X_3)\geq 1$ for a
nonzero range, the displayed factored estimate follows.
\end{proof}

For completeness, attainment follows by fixing a basis of the target, writing
a bounded minimizing net of projections in terms of coordinate functionals,
and taking weak-star convergent subnets of those functionals.  The projection
identities pass to the limit and the operator norm is lower semicontinuous.

\begin{corollary}[Fixed-dimensional upper bound]\label{cor:upper}
If $\dim X_1=d$, then the conclusion requested in Problem~3 holds for that
triple with
\[
        C_d=1+2\sqrt d.
\]
\end{corollary}

\begin{proof}
The Kadets--Snobar theorem \cite{KadetsSnobar} gives $\lambda(X_1,X_3)\leq\sqrt d$; apply
Proposition~\ref{prop:correction}.
\end{proof}

\section{An obstruction from a unique minimal projection}
\label{sec:obstruction}

The key observation is abstract.  A unique minimal projection with large norm
forces a bad intermediate hyperplane.

\begin{lemma}\label{lem:quotient}
Let $X$ be finite-dimensional, let $E\subsetneq X$, and let $P:X\to E$ be a
projection.  Put $K=\ker P$, let $q:X\to X/E$ be the quotient map, and define
\[
       J=\restr{q}{K}:K\longrightarrow X/E.
\]
Then $J$ is an isomorphism and
\[
       \|J^{-1}\|=\|\Id-P\|.
\]
Consequently, there is a functional $\varphi\in E^\perp$ such that
\begin{equation}\label{eq:phi}
       \|\varphi\|=1,
       \qquad
       \bigl\|\restr{\varphi}{K}\bigr\|
       =\frac{1}{\|\Id-P\|}.
\end{equation}
\end{lemma}

\begin{proof}
Because $X=E\oplus K$, the map $J$ is bijective and
\[
       J^{-1}(q x)=(\Id-P)x.
\]
For every $x\in X$ and $e\in E$,
$(\Id-P)(x-e)=(\Id-P)x$.  Taking an infimum over $e$ gives
\[
 \|(\Id-P)x\|\leq \|\Id-P\|\,\|qx\|,
\]
so $\|J^{-1}\|\leq\|\Id-P\|$.  Conversely,
\[
 \|(\Id-P)x\|=\|J^{-1}(qx)\|
 \leq \|J^{-1}\|\,\|qx\|
 \leq \|J^{-1}\|\,\|x\|,
\]
and equality of the two norms follows.

Since $J^*$ is an isomorphism between finite-dimensional spaces,
\[
 \min_{\|\psi\|=1}\|J^*\psi\|
   =\frac{1}{\|(J^*)^{-1}\|}
   =\frac{1}{\|J^{-1}\|}.
\]
Choose a unit vector $\psi\in(X/E)^*$ attaining this minimum and set
$\varphi=q^*\psi$.  The adjoint $q^*$ is an isometry from $(X/E)^*$ onto
$E^\perp$, and $\restr{\varphi}{K}=J^*\psi$.  This proves
\eqref{eq:phi}.
\end{proof}

\begin{proposition}[Bad intermediate hyperplane]\label{prop:hyperplane}
Let $X$ be finite-dimensional, let $E\subsetneq X$, and suppose that
$P:X\to E$ is the unique minimal projection.  There is a hyperplane
$F\subset X$ containing $E$ such that
\[
       \lambda(F,X)\leq 2
\]
and the following holds: whenever $A:F\to E$ and $R:X\to F$ are projections
for which $AR$ is minimal onto $E$, one has
\begin{equation}\label{eq:badR}
       \|R\|\geq \|\Id-P\|-1\geq \|P\|-2.
\end{equation}
\end{proposition}

\begin{proof}
Let $K=\ker P$, choose $\varphi$ as in Lemma~\ref{lem:quotient}, and set
$F=\ker\varphi$.  Since $\varphi\in E^\perp$, we have $E\subset F$.
A norm-one functional on a finite-dimensional space attains its norm, so one
can choose $v\in X$ with $\|v\|=1$ and $\varphi(v)=1$.  Then
\[
       R_0=\Id-v\otimes\varphi
\]
is a projection onto $F$ and $\|R_0\|\leq 2$.  Hence
$\lambda(F,X)\leq2$.

Now let $A$ and $R$ be as in the statement.  The product $AR$ is a projection
onto $E$.  Since it is minimal and $P$ is the unique minimal projection,
$AR=P$.  Every projection onto the hyperplane $F=\ker\varphi$ has the form
\[
       R=\Id-u\otimes\varphi,
       \qquad \varphi(u)=1.
\]
Because $\ker R\subset\ker(AR)=\ker P=K$, the vector $u$ belongs to $K$.
Thus \eqref{eq:phi} yields
\[
  1=|\varphi(u)|
    \leq \bigl\|\restr{\varphi}{K}\bigr\|\,\|u\|
    =\frac{\|u\|}{\|\Id-P\|},
\]
so $\|u\|\geq\|\Id-P\|$.  Since $\|\varphi\|=1$,
\[
 \|\Id-R\|=\|u\otimes\varphi\|=\|u\|\geq\|\Id-P\|.
\]
The triangle inequality gives $\|R\|\geq\|\Id-R\|-1$, proving the first
inequality in \eqref{eq:badR}.  The second follows from
$\|P\|\leq1+\|\Id-P\|$.
\end{proof}

\section{Unique minimal Rademacher projections of growing norm}
\label{sec:rademacher}

Let
\[
 \Omega_d=\{-1,1\}^d
\]
with uniform probability measure, let $p_d=2d$, and put
\[
 X_d=L_{p_d}(\Omega_d).
\]
For $1\leq j\leq d$, write $r_j(\omega)=\omega_j$ and set
\[
 E_d=\spann\{r_1,\ldots,r_d\}.
\]
Define the first-chaos, or Rademacher, projection by
\begin{equation}\label{eq:Rd}
 \mathcal R_d f
   =\sum_{j=1}^d \E(f r_j)\,r_j.
\end{equation}

\begin{lemma}[Unique minimality]\label{lem:unique}
For every $d\geq2$, the operator $\mathcal R_d:X_d\to E_d$ is the unique
minimal projection onto $E_d$.
\end{lemma}

\begin{proof}
Identify $X_d$ with the subspace of $L_{p_d}[0,1]$ consisting of functions
measurable with respect to the sigma-algebra generated by the first $d$
classical Rademacher functions.  Let $L_d$ be conditional expectation onto
this finite sigma-algebra; then $\|L_d\|=1$.  Let
$\widetilde{\mathcal R}_d$ be the classical Rademacher projection on
$L_{p_d}[0,1]$, and let $T_d$ be its restriction to $X_d$.  Since the
Rademacher functions are measurable with respect to the finite sigma-algebra,
\[
       \widetilde{\mathcal R}_d=T_dL_d.
\]
Lewicki and Skrzypek proved that $\widetilde{\mathcal R}_d$ is the unique
minimal projection onto $E_d$ in $L_p[0,1]$ for every $1<p<\infty$
\cite{LewickiSkrzypek,Skrzypek}.

We transfer that statement to $X_d$.  Restriction gives
$\|T_d\|\leq\|\widetilde{\mathcal R}_d\|$, while the displayed factorization
and $\|L_d\|=1$ give the reverse inequality.  If $Q:X_d\to E_d$ is any
projection, then $QL_d$ is a projection from $L_{p_d}[0,1]$ onto $E_d$.
It follows that
\[
 \|T_d\|=\lambda(E_d,L_{p_d}[0,1])\leq\|Q\|.
\]
Thus $T_d$ is minimal on $X_d$.  If $Q$ is another minimal projection on
$X_d$, then $QL_d$ is minimal on $L_{p_d}[0,1]$ and hence, by uniqueness,
$QL_d=\widetilde{\mathcal R}_d$.  Restriction to $X_d$ gives $Q=T_d$.
Under the finite-cube identification, $T_d$ is exactly \eqref{eq:Rd}.
\end{proof}

\begin{lemma}[Norm growth]\label{lem:growth}
For every $d\geq2$,
\[
       \|\mathcal R_d\|_{X_d\to X_d}\geq \sqrt{d/6}.
\]
\end{lemma}

\begin{proof}
The same linear map acts on every $L_p(\Omega_d)$.  Its kernel formula is
\[
 (\mathcal R_df)(\omega)
   =2^{-d}\sum_{\eta\in\Omega_d}f(\eta)
       \sum_{j=1}^d\omega_j\eta_j.
\]
Therefore its $L_\infty$ operator norm is the common absolute row sum
\[
 \|\mathcal R_d\|_{\infty\to\infty}
   =\beta_d
   :=\E\left|\sum_{j=1}^d\varepsilon_j\right|,
\]
where the $\varepsilon_j$ are independent symmetric signs.

On a probability space with $N$ atoms of equal mass,
\[
       N^{-1/p}\|g\|_\infty\leq\|g\|_p\leq\|g\|_\infty.
\]
Choosing an $L_\infty$ norming vector for the finite matrix of
$\mathcal R_d$ gives
\[
 \|\mathcal R_d\|_{p\to p}
     \geq N^{-1/p}\|\mathcal R_d\|_{\infty\to\infty}.
\]
Here $N=2^d$ and $p=p_d=2d$, so
\begin{equation}\label{eq:lp-linf}
       \|\mathcal R_d\|_{p_d\to p_d}\geq 2^{-1/2}\beta_d.
\end{equation}

Let $S_d=\sum_{j=1}^d\varepsilon_j$.  H\"older's inequality gives
\[
 \E S_d^2
 \leq (\E|S_d|)^{2/3}(\E S_d^4)^{1/3},
\]
and hence
\[
 \E|S_d|\geq
 \frac{(\E S_d^2)^{3/2}}{(\E S_d^4)^{1/2}}.
\]
Since
\[
 \E S_d^2=d,
 \qquad
 \E S_d^4=3d^2-2d\leq3d^2,
\]
we obtain $\beta_d\geq\sqrt{d/3}$.  Combining this with
\eqref{eq:lp-linf} proves the lemma.
\end{proof}

\begin{proof}[Proof of Theorem~\ref{thm:main}]
Apply Proposition~\ref{prop:hyperplane} to the unique minimal projection
$\mathcal R_d:X_d\to E_d$.  It supplies a hyperplane $F_d$ containing
$E_d$ such that $\lambda(F_d,X_d)\leq2$.  Set
\[
       X_{1,d}=E_d,\qquad X_{2,d}=F_d,\qquad X_{3,d}=X_d.
\]
If $P_1:F_d\to E_d$ and $P_2:X_d\to F_d$ are projections with
$\|P_1P_2\|=\lambda(E_d,X_d)$, then $P_1P_2$ is a minimal projection onto
$E_d$.  Lemma~\ref{lem:unique} makes it equal to $\mathcal R_d$, and
Proposition~\ref{prop:hyperplane}, followed by Lemma~\ref{lem:growth}, gives
\[
       \|P_2\|\geq\|\mathcal R_d\|-2
                 \geq\sqrt{d/6}-2.
\]
Since $\lambda(F_d,X_d)\leq2$, the ratio in the theorem tends to infinity.
\end{proof}

\section{Consequences and literature status}

Define $C_d^*$ to be the least constant that works uniformly in Problem~3 over
all triples with $\dim X_1=d$.  Corollary~\ref{cor:upper} and
Theorem~\ref{thm:main} show
\[
 \frac{\sqrt{d/6}-2}{2}\leq C_d^*\leq1+2\sqrt d.
\]
Thus the $\sqrt d$ dependence in the supplied partial result is optimal in
order, although the numerical constants have not been optimized.

A targeted search of the title and exact wording of Ostrovskii's Problem~3,
together with searches through visible citation trails and projection-constant
literature, did not locate a paper explicitly giving a full answer or
counterexample.  This is not an exhaustive priority claim.  The construction
above uses a published result that predates the problem: the unique minimality
of Rademacher projections proved by Lewicki and Skrzypek and restated explicitly
by Skrzypek.  All other ingredients in the counterexample are included here.

\section*{Human review notes}

The finite-rank correction formula in Proposition~\ref{prop:correction} is
elementary and is retained because it gives the matching upper estimate.  The
main mathematical review point for the negative answer is
Lemma~\ref{lem:unique}: one should check the cited unique minimality theorem
for the classical Rademacher projection and the conditional-expectation
transfer from $L_{2d}[0,1]$ to the finite cube.  Once that input is accepted,
the hyperplane obstruction and norm-growth estimates are finite-dimensional
and self-contained.

\begin{thebibliography}{9}

\bibitem{KadetsSnobar}
M.~I. Kadets and M.~G. Snobar,
\emph{Certain functionals on the Minkowski compactum},
Math. Notes \textbf{10} (1971), 694--696.

\bibitem{LewickiSkrzypek}
G. Lewicki and L. Skrzypek,
\emph{Chalmers--Metcalf operator and uniqueness of minimal projections},
J. Approx. Theory \textbf{148} (2007), no.~1, 71--91.
\href{https://doi.org/10.1016/j.jat.2006.10.008}{doi:10.1016/j.jat.2006.10.008}.

\bibitem{Ostrovskii}
M.~I. Ostrovskii,
\emph{Compositions of projections in Banach spaces and relations between
approximation properties},
Rocky Mountain J. Math. \textbf{38} (2008), no.~4, 1253--1262.
\href{https://doi.org/10.1216/RMJ-2008-38-4-1253}{doi:10.1216/RMJ-2008-38-4-1253}.

\bibitem{Skrzypek}
L. Skrzypek,
\emph{On the $L_p$ norm of the Rademacher projection and related inequalities},
Proc. Amer. Math. Soc. \textbf{137} (2009), no.~8, 2661--2669.
\href{https://doi.org/10.1090/S0002-9939-09-09875-X}{doi:10.1090/S0002-9939-09-09875-X}.

\end{thebibliography}

\end{document}
